\documentclass{amsart}
\usepackage{amsmath,amssymb,amsthm,hyperref}
\usepackage{algorithm}
\usepackage{algpseudocode}
\usepackage{bm}

\newtheorem{theorem}{Theorem}
\newtheorem{lemma}{Lemma}
\newtheorem{proposition}{Proposition}
\newtheorem{corollary}{Corollary}
\theoremstyle{definition}
\newtheorem{definition}{Definition}
\newtheorem{remark}{Remark}
\newtheorem{assumption}{Assumption}

\DeclareMathOperator{\spec}{spec}
\DeclareMathOperator{\diag}{diag}
\DeclareMathOperator{\dist}{dist}
\newcommand{\R}{\mathbb{R}}
\newcommand{\C}{\mathbb{C}}
\newcommand{\Z}{\mathbb{Z}}
\newcommand{\norm}[1]{\left\lVert#1\right\rVert}
\newcommand{\Acal}{\mathcal{A}}
\newcommand{\Bcal}{\mathcal{B}}
\newcommand{\Ccal}{\mathcal{C}}
\newcommand{\Dcal}{\mathcal{D}}
\newcommand{\Hcal}{\mathcal{H}}
\newcommand{\Ncal}{\mathcal{N}}

\begin{document}

\title{Numerically reliable, structure-exploiting computation of the frequency
response of linear time-periodic systems}
\maketitle

\section{Formalization of the problem}

\subsection{The system and the standing assumptions}
We are given a continuous-time linear time-periodic (LTP) state-space system
\begin{equation}\label{eq:sys}
\dot x(t)=A(t)x(t)+B(t)u(t),\qquad y(t)=C(t)x(t)+D(t)u(t),
\end{equation}
with $x(t)\in\C^{n}$, $u(t)\in\C^{m}$, $y(t)\in\C^{p}$, and with $A,B,C,D$
matrix-valued, $T$-periodic, $A(t+T)=A(t)$ etc., $T>0$. Set $\Omega:=2\pi/T$.
We assume the data are bounded measurable on $[0,T)$; this is the case for all
standard representations (truncated Fourier series, piecewise-constant,
band-limited sampled coefficients). We write the Fourier series
\begin{equation}\label{eq:fourier}
A(t)=\sum_{k\in\Z}\hat A_k\,e^{ik\Omega t},\qquad
\hat A_k=\frac1T\int_0^T A(t)e^{-ik\Omega t}\,dt,
\end{equation}
and analogously $\hat B_k,\hat C_k,\hat D_k$. Reality of the data is equivalent
to $\hat A_{-k}=\overline{\hat A_k}$, etc.; we never use reality below, so the
results hold verbatim for complex periodic data.

\begin{assumption}[Exponential dichotomy / hyperbolicity on the evaluation line]
\label{ass:hyp}
Let $\Phi_A(t,s)$ denote the state-transition matrix of $\dot z=A(t)z$, i.e.
$\partial_t\Phi_A(t,s)=A(t)\Phi_A(t,s)$, $\Phi_A(s,s)=I$, and let
$\Psi:=\Phi_A(T,0)$ be the \emph{monodromy matrix}. The \emph{Floquet
characteristic multipliers} are the eigenvalues $\mu_1,\dots,\mu_n$ of $\Psi$;
fixing a branch of the logarithm, the \emph{Floquet exponents} are
$\lambda_j:=\tfrac1T\log\mu_j$ (defined modulo $i\Omega$). For every frequency
$\omega$ at which we evaluate, we assume
\begin{equation}\label{eq:hyp}
i\omega\notin\spec(\Psi^{1/T})+i\Omega\Z,\qquad\text{equivalently}\qquad
e^{i\omega T}\notin\spec(\Psi).
\end{equation}
\end{assumption}

Condition \eqref{eq:hyp} is exactly the condition that the boundary-value
problem appearing below is uniquely solvable; it holds for all real $\omega$
when the system is exponentially stable (all $\operatorname{Re}\lambda_j<0$), in
which case it holds for all real $\omega$ \emph{simultaneously}. The set of
$\omega$ violating \eqref{eq:hyp} is the (discrete, $\Omega$-periodic) pole set
of the frequency response. Two distinct distance-to-singularity
quantities will appear, and it is essential to keep them apart because they
govern \emph{different} numerical objects.
\begin{itemize}
\item The \emph{generator gap}, which controls the norm of the (lifted) resolvent
$\mathcal G(\omega)$ and hence the off-diagonal decay rate used by the
truncation algorithm:
\begin{equation}\label{eq:gengap}
\gamma(\omega):=\dist\!\Big(i\omega,\ \big\{\lambda_j-ik\Omega:\ j=1,\dots,n,\ k\in\Z\big\}\Big)
=\min_{j,k}\big|i\omega-\lambda_j+ik\Omega\big|>0 .
\end{equation}
\item The \emph{monodromy gap}, which controls the conditioning of the
finite two-point BVP solve in the monodromy/Floquet algorithms:
\begin{equation}\label{eq:gap}
\delta(\omega):=\dist\!\big(1,\ \spec(e^{-i\omega T}\Psi)\big)
=\dist\!\big(e^{i\omega T},\spec(\Psi)\big)>0 .
\end{equation}
\end{itemize}
Both are positive precisely under Assumption~\ref{ass:hyp}, and both vanish as
$\omega$ approaches a pole. \emph{Warning.} It is \emph{not} true in general that
$\|\mathcal G(\omega)\|\lesssim 1/\delta(\omega)$: when the Floquet multipliers
are large (unstable modes), $\delta(\omega)$ can be huge while
$\|\mathcal G(\omega)\|$ stays $O(1)$. The correct resolvent bound uses
$\gamma(\omega)$ (Proposition~\ref{prop:resnorm}). Conversely the BVP-solve
conditioning is genuinely $\Theta(1/\delta(\omega))$ up to eigenvector
conditioning. We use each gap only where it is provably the right quantity.

\subsection{The harmonic transfer function (EMP frequency-response operator)}
\label{sec:HTF}
An \emph{exponentially modulated periodic} (EMP) input at fundamental frequency
$\omega$ is $u(t)=\sum_{k\in\Z}u_k\,e^{i(\omega+k\Omega)t}$. Substituting the
EMP ansatz $x(t)=\sum_k x_k e^{i(\omega+k\Omega)t}$,
$y(t)=\sum_k y_k e^{i(\omega+k\Omega)t}$ into \eqref{eq:sys} and matching the
$e^{i(\omega+k\Omega)t}$ harmonic (using
$A(t)x(t)=\sum_k\big(\sum_\ell\hat A_{k-\ell}x_\ell\big)e^{i(\omega+k\Omega)t}$)
yields the doubly-infinite algebraic relations
\begin{equation}\label{eq:harmbal}
i(\omega+k\Omega)x_k=\sum_{\ell}\hat A_{k-\ell}x_\ell+\sum_\ell\hat B_{k-\ell}u_\ell,
\quad
y_k=\sum_\ell\hat C_{k-\ell}x_\ell+\sum_\ell\hat D_{k-\ell}u_\ell .
\end{equation}
Collecting the harmonics into bi-infinite vectors
$\mathbf{x}=(x_k)_{k\in\Z}$, $\mathbf u,\mathbf y$ similarly, define the
\emph{block-Laurent (Toeplitz) operators}
\begin{equation}\label{eq:Toep}
(\Acal)_{k\ell}=\hat A_{k-\ell},\quad
(\Bcal)_{k\ell}=\hat B_{k-\ell},\quad
(\Ccal)_{k\ell}=\hat C_{k-\ell},\quad
(\Dcal)_{k\ell}=\hat D_{k-\ell},
\end{equation}
acting on $\ell^2(\Z;\C^n)$ (resp. $\C^m,\C^p$), and the diagonal modulation
operator $\Ncal:=\diag(k I_n)_{k\in\Z}$. Then \eqref{eq:harmbal} reads
$\big(i(\omega I+\Omega\Ncal)-\Acal\big)\mathbf x=\Bcal\mathbf u$ and
$\mathbf y=\Ccal\mathbf x+\Dcal\mathbf u$, so the \emph{harmonic transfer
function} (HTF), the frequency-response operator, is
\begin{equation}\label{eq:HTFdef}
\boxed{\;\Hcal(\omega)=\Ccal\big(i(\omega I+\Omega\Ncal)-\Acal\big)^{-1}\Bcal+\Dcal\;}
\end{equation}
acting on $\ell^2(\Z;\C^m)\to\ell^2(\Z;\C^p)$, with $\mathbf y=\Hcal(\omega)\mathbf u$.
The HTF is $\Omega$-quasiperiodic in $\omega$:
$\Hcal(\omega+\Omega)=S^*\Hcal(\omega)S$ where $S$ is the bilateral shift, so it
suffices to compute it on a \emph{fundamental band} $\omega\in[0,\Omega)$ (or
$(-\Omega/2,\Omega/2]$). The $(k,\ell)$ block $\Hcal(\omega)_{k\ell}\in\C^{p\times m}$
maps the input harmonic $u_\ell e^{i(\omega+\ell\Omega)t}$ to its contribution
to the output harmonic $y_k e^{i(\omega+k\Omega)t}$.

\medskip
We were asked to: (i) compute $\Hcal(\omega)$ on a grid of frequencies, in a
numerically reliable way, with provable error bounds and favorable complexity;
(ii) avoid forming/decomposing the doubly-infinite operator at each frequency;
(iii) exploit structure (sparsity / harmonic roll-off / monodromy–Floquet
factorizations); (iv) match the efficiency available in the LTI case.

Sections~\ref{sec:wellposed}--\ref{sec:trunc} establish well-posedness and a
truncation algorithm with a sharp error bound. Section~\ref{sec:mono} gives the
monodromy/Floquet algorithm that avoids the infinite operator entirely and is
the recommended method; Section~\ref{sec:lift} gives the
fast-frequency-sweep refinement. Section~\ref{sec:complexity} collects the
complexity statements.

\section{Well-posedness of the operator definition}\label{sec:wellposed}

We must first show \eqref{eq:HTFdef} is meaningful: that
$i(\omega I+\Omega\Ncal)-\Acal$ is boundedly invertible under
Assumption~\ref{ass:hyp}, and identify its inverse with a Green's operator that
we can compute in finite dimensions.

\begin{lemma}[Floquet operator conjugation]\label{lem:conj}
Let $P(t)$ be the $T$-periodic Lyapunov--Floquet factor of $A$, with
$P,P^{-1}\in L^\infty\cap W^{1,\infty}$ (true for all standard data; $P$ is
absolutely continuous because $\Phi_A$ is, and $P,P^{-1}$ are bounded by
periodicity and continuity), and let $R$ be the constant matrix with
$\Phi_A(t,0)=P(t)e^{Rt}$, $P(0)=I$, $e^{RT}=\Psi$. Let
$\mathcal P=(\hat P_{k-\ell})_{k\ell}$ and $\mathcal P'=(\widehat{P'}_{k-\ell})_{k\ell}$
be the block-Laurent operators of $P$ and of its derivative $\dot P=P'$,
$\mathcal R=\diag(R)_{k\in\Z}$ the constant block-diagonal operator, and
$\Ncal=\diag(kI_n)_{k}$. Then $\mathcal P,\mathcal P^{-1}=(\widehat{P^{-1}}_{k-\ell})$
are bounded with bounded inverses on $\ell^2(\Z;\C^n)$, and
\begin{equation}\label{eq:commNP}
[\,i\Omega\Ncal,\ \mathcal P\,]\;=\;\mathcal P',
\end{equation}
\begin{equation}\label{eq:conj}
\mathcal P^{-1}\big(i\Omega\Ncal-\Acal\big)\mathcal P
\;=\;i\Omega\Ncal-\mathcal R .
\end{equation}
\end{lemma}

\begin{proof}
Boundedness of $\mathcal P,\mathcal P^{-1}$ with the stated norms is
Lemma~\ref{lem:bounded} applied to $P$ and $P^{-1}$; that $\mathcal P^{-1}$ (the
operator inverse) coincides with the Laurent operator of the pointwise inverse
$P(t)^{-1}$ is the statement that the multiplication operators $M_P,M_{P^{-1}}$
on $L^2([0,T];\C^n)$ are mutually inverse, which is immediate since
$M_P M_{P^{-1}}=M_{PP^{-1}}=M_I=I$, and the Laurent matrix is the matrix of the
multiplication operator in the Fourier basis (Lemma~\ref{lem:bounded}).

For \eqref{eq:commNP}: the operator $i\Omega\Ncal$ is, in the Fourier basis, the
matrix of the differentiation operator $\tfrac{d}{dt}$ on $T$-periodic functions,
since $\tfrac{d}{dt}e^{ik\Omega t}=ik\Omega\,e^{ik\Omega t}$. The commutator of
$\tfrac{d}{dt}$ with the multiplication operator $M_P$ is, by the Leibniz rule,
$[\tfrac{d}{dt},M_P]f=\tfrac{d}{dt}(Pf)-P\tfrac{d}{dt}f=\dot P\,f=M_{\dot P}f$,
i.e. $[\,d/dt,\,M_P]=M_{\dot P}$. Passing to the Fourier matrix representation
(an isometry of $L^2$ onto $\ell^2$ that carries $M_P\mapsto\mathcal P$,
$M_{\dot P}\mapsto\mathcal P'$, and $\tfrac{d}{dt}\mapsto i\Omega\Ncal$) yields
\eqref{eq:commNP}. Entrywise this is the identity
$(i\Omega k)\hat P_{k-\ell}-\hat P_{k-\ell}(i\Omega\ell)=i\Omega(k-\ell)\hat P_{k-\ell}
=\widehat{\dot P}_{k-\ell}$, the last equality being the Fourier rule
$\widehat{\dot P}_{m}=i\Omega m\,\hat P_m$; this also shows $\mathcal P'$ is
bounded since $\dot P\in L^\infty$.

For \eqref{eq:conj}: rearrange \eqref{eq:commNP} as
$(i\Omega\Ncal)\mathcal P=\mathcal P(i\Omega\Ncal)+\mathcal P'$, so
\[
\mathcal P^{-1}(i\Omega\Ncal)\mathcal P
=i\Omega\Ncal+\mathcal P^{-1}\mathcal P' .
\]
Hence
\[
\mathcal P^{-1}(i\Omega\Ncal-\Acal)\mathcal P
=i\Omega\Ncal+\mathcal P^{-1}\mathcal P'-\mathcal P^{-1}\Acal\mathcal P
=i\Omega\Ncal-\mathcal P^{-1}\big(\Acal\mathcal P-\mathcal P'\big).
\]
It remains to identify $\mathcal P^{-1}(\Acal\mathcal P-\mathcal P')$ with the
constant block-diagonal operator $\mathcal R$. Recall (Lemma~\ref{lem:bounded})
that the map $F\mapsto(\hat F_{k-\ell})_{k\ell}$ from $L^\infty([0,T])$ matrix
symbols to block-Laurent operators is the Fourier matrix representation of the
multiplication operator $M_F$; since $M_F M_G=M_{FG}$, this map is an algebra
homomorphism, so products and differences of Laurent operators are again Laurent
operators with the product/difference symbol. In particular $\Acal\mathcal P$ has
symbol $A(t)P(t)$, $\Acal\mathcal P-\mathcal P'$ has symbol $A(t)P(t)-\dot P(t)$,
and $\mathcal P^{-1}$ has symbol $P(t)^{-1}$. Thus it suffices to identify the
symbol of $\mathcal P^{-1}(\Acal\mathcal P-\mathcal P')$: the
Laurent operator of $\Acal\mathcal P-\mathcal P'$ is multiplication by the
$T$-periodic matrix $A(t)P(t)-\dot P(t)$, and $\mathcal P^{-1}$ is multiplication
by $P(t)^{-1}$, so $\mathcal P^{-1}(\Acal\mathcal P-\mathcal P')$ is multiplication
by $P(t)^{-1}\big(A(t)P(t)-\dot P(t)\big)$. But differentiating
$\Phi_A(t,0)=P(t)e^{Rt}$ gives
$A(t)P(t)e^{Rt}=\dot\Phi_A=\dot P(t)e^{Rt}+P(t)Re^{Rt}$, i.e.
$A(t)P(t)-\dot P(t)=P(t)R$ (multiply on the right by $e^{-Rt}$), hence
$P(t)^{-1}\big(A(t)P(t)-\dot P(t)\big)=R$, a \emph{constant} matrix. The Laurent
operator of the constant symbol $R$ is exactly $\mathcal R=\diag(R)_k$. This
proves \eqref{eq:conj}.
\end{proof}

\begin{lemma}[Spectrum of the lifted generator]\label{lem:spectrum}
Let $\Acal$ be the bounded block-Laurent operator \eqref{eq:Toep} (bounded
because $A\in L^\infty$, see Lemma~\ref{lem:bounded}). Then
\begin{equation}\label{eq:specG}
\spec\!\big(i\Omega\Ncal-\Acal+i\omega I\big)
=\big\{\,i\omega-\lambda_j+i k\Omega:\ j=1,\dots,n,\ k\in\Z\,\big\}^{\mathrm{cl}},
\end{equation}
where $\lambda_j$ are the Floquet exponents and the closure adds the limit
points; in particular $i(\omega I+\Omega\Ncal)-\Acal$ is invertible iff
$i\omega\ne\lambda_j-ik\Omega$ for all $j,k$, which is exactly \eqref{eq:hyp}.
\end{lemma}

\begin{proof}
Consider the operator $\mathsf G:=i\Omega\Ncal-\Acal$ as the generator
representation of the LTP flow. By Floquet theory there exists a $T$-periodic,
boundedly invertible $P(t)$ (the Lyapunov–Floquet transformation) and a
constant matrix $R$ with $\spec(R)=\{\lambda_1,\dots,\lambda_n\}$ such that
$\Phi_A(t,0)=P(t)e^{Rt}$, $P(0)=I$, $e^{RT}=\Psi$. Equivalently, the change of
variables $z=P(t)\zeta$ transforms $\dot z=A(t)z$ into the LTI system
$\dot\zeta=R\zeta$. Passing to harmonics, Lemma~\ref{lem:conj} gives the exact
operator identity $\mathcal P^{-1}\mathsf G\,\mathcal P=i\Omega\Ncal-\mathcal R$,
where $\mathcal P=(\hat P_{k-\ell})$ is bounded with bounded inverse and
$\mathcal R=\diag(R)_{k\in\Z}$ is block-diagonal (the constant matrix $R$ has
only the $k=0$ Fourier coefficient). The nontrivial point---that $\Ncal$ does
\emph{not} commute with $\mathcal P$, yet the conjugation still produces the
\emph{constant}-coefficient generator---is exactly \eqref{eq:conj}, whose proof
uses the commutator identity $[\,i\Omega\Ncal,\mathcal P\,]=\mathcal P'$ together
with the Floquet ODE relation $A(t)P(t)-\dot P(t)=P(t)R$. The operator $i\Omega\Ncal-\mathcal R$ is
block-diagonal in $k$ with $k$-th block $ik\Omega I-R$; its spectrum is
$\bigcup_k\spec(ik\Omega I-R)=\{ik\Omega-\lambda_j\}$, with closure adding the
accumulation at $\infty$ along $i\Omega\Z$. Since similarity preserves spectrum
and the spectrum of $\mathsf G+i\omega I$ is that of $\mathsf G$ shifted by
$i\omega$, \eqref{eq:specG} follows. Invertibility of
$i\omega I-\mathsf G$ is the statement that $0\notin\spec(i\omega I-\mathsf G)$,
i.e. $i\omega\ne\lambda_j-ik\Omega$; exponentiating, $e^{i\omega T}\ne
e^{\lambda_j T}=\mu_j$, which is \eqref{eq:hyp}.
\end{proof}

\begin{lemma}[Boundedness of the Toeplitz operators]\label{lem:bounded}
If $A\in L^\infty([0,T];\C^{n\times n})$ then $\Acal$ in \eqref{eq:Toep} is a
bounded operator on $\ell^2(\Z;\C^n)$ with
$\norm{\Acal}_{\ell^2\to\ell^2}=\operatorname*{ess\,sup}_{t}\norm{A(t)}_2
=:\norm{A}_{L^\infty}$. The same holds for $\Bcal,\Ccal,\Dcal$.
\end{lemma}

\begin{proof}
$\Acal$ is the matrix of the multiplication operator
$M_A:f\mapsto A(\cdot)f(\cdot)$ on $L^2([0,T];\C^n)$ in the orthonormal Fourier
basis $\{e^{ik\Omega t}\}$. Indeed, if $f=\sum_\ell f_\ell e^{i\ell\Omega t}$
then $Af=\sum_k(\sum_\ell\hat A_{k-\ell}f_\ell)e^{ik\Omega t}$, so the matrix of
$M_A$ is precisely $(\hat A_{k-\ell})_{k\ell}=\Acal$. The norm of a
multiplication operator on $L^2$ equals the essential supremum of the pointwise
operator norm, $\norm{M_A}=\operatorname*{ess\,sup}_t\norm{A(t)}_2$. Since the
Fourier basis is orthonormal, the operator and its $\ell^2$ matrix
representation have equal norm.
\end{proof}

\begin{proposition}[Norm of the lifted resolvent via the generator gap]
\label{prop:resnorm}
Under Assumption~\ref{ass:hyp} the resolvent
$\mathcal G(\omega)=\big(i(\omega I+\Omega\Ncal)-\Acal\big)^{-1}$ is bounded with
\begin{equation}\label{eq:resnorm}
\norm{\mathcal G(\omega)}_{\ell^2\to\ell^2}\le \frac{\kappa_\star}{\gamma(\omega)},
\qquad
\kappa_\star:=\kappa(\mathcal P)\,\kappa(V),\quad
\kappa(\mathcal P):=\norm{\mathcal P}\,\norm{\mathcal P^{-1}}
=\Big(\operatorname*{ess\,sup}_t\norm{P(t)}_2\Big)\Big(\operatorname*{ess\,sup}_t\norm{P(t)^{-1}}_2\Big),
\end{equation}
where $P$ is the Lyapunov--Floquet transformation of Lemma~\ref{lem:spectrum},
$\gamma(\omega)$ is the generator gap \eqref{eq:gengap}, and $\kappa(V)=\norm
V\norm{V^{-1}}$ is the eigenvector conditioning of the constant matrix $R$
(equivalently of the monodromy $\Psi$), with $\kappa(V)=1$ when $\Psi$ is normal.
Conversely $\norm{\mathcal G(\omega)}\ge 1/\big(\gamma(\omega)\,\kappa_\star\big)$,
so up to the data-dependent but $\omega$-independent constant $\kappa_\star$ the
resolvent norm is $\Theta(1/\gamma(\omega))$. (In the non-diagonalizable case
$\kappa(V)$ is replaced by the corresponding Jordan-block resolvent constant; the
bound remains $\Theta(1/\gamma(\omega))$ up to an $\omega$-independent factor and
a possible polynomial-in-$1/\gamma$ correction of degree below the largest Jordan
block.) In general
$\norm{\mathcal G(\omega)}$ is \emph{not} controlled by $1/\delta(\omega)$.
\end{proposition}

\begin{proof}
By Lemma~\ref{lem:conj} (with the trivial added scalar $i\omega I$, which
commutes with everything), conjugation by the bounded, boundedly invertible
block-Laurent operator $\mathcal P=(\hat P_{k-\ell})$ gives the exact identity
$\mathcal P^{-1}\big(i(\omega I+\Omega\Ncal)-\Acal\big)\mathcal P
=i(\omega I+\Omega\Ncal)-\mathcal R$, where $\mathcal R=\diag(R)_k$ is
block-diagonal with the constant matrix $R$ ($\spec R=\{\lambda_j\}$) in every
block. Here $\Ncal$ does \emph{not} commute with $\mathcal P$; rather the
commutator $[\,i\Omega\Ncal,\mathcal P\,]=\mathcal P'$ \eqref{eq:commNP} supplies
precisely the $\dot P$ term that turns the conjugated generator into the
constant-coefficient one, as proved in Lemma~\ref{lem:conj}. Hence
\[
\mathcal G(\omega)=\mathcal P\,\big(i(\omega I+\Omega\Ncal)-\mathcal R\big)^{-1}\,\mathcal P^{-1},
\]
and the middle factor is block-diagonal in $k$ with $k$-th block
$\big(i(\omega+k\Omega)I-R\big)^{-1}$. Therefore
\[
\norm{\big(i(\omega I+\Omega\Ncal)-\mathcal R\big)^{-1}}
=\sup_k\norm{\big(i(\omega+k\Omega)I-R\big)^{-1}}_2 .
\]
Now $R=V\,\diag(\lambda_j)\,V^{-1}$ (we may take $R$ in this diagonalizable form
when $\Psi$ is diagonalizable; the general Jordan case adds only the standard
finite eigenvector-conditioning factor and replaces the bound by the resolvent
of a Jordan block, still $\Theta(1/\gamma)$). For a fixed $k$, using $R=V\diag(\lambda_j)V^{-1}$ and
$\norm{(i(\omega+k\Omega)I-\diag(\lambda_j))^{-1}}_2
=1/\min_j|i(\omega+k\Omega)-\lambda_j|$,
\[
\norm{(i(\omega+k\Omega)I-R)^{-1}}_2
\le \frac{\kappa(V)}{\min_j|i(\omega+k\Omega)-\lambda_j|}
= \frac{\kappa(V)}{\min_j|i\omega-\lambda_j+ik\Omega|}.
\]
Taking the supremum over $k\in\Z$,
\[
\sup_k\norm{(i(\omega+k\Omega)I-R)^{-1}}_2
\le \frac{\kappa(V)}{\min_{j,k}|i\omega-\lambda_j+ik\Omega|}
= \frac{\kappa(V)}{\gamma(\omega)}.
\]
Multiplying by $\norm{\mathcal P}\norm{\mathcal P^{-1}}=\kappa(\mathcal P)$ from
the two outer factors gives
$\norm{\mathcal G(\omega)}\le\kappa(\mathcal P)\kappa(V)/\gamma(\omega)
=\kappa_\star/\gamma(\omega)$, which is \eqref{eq:resnorm}. For the lower bound, write the middle factor
$\mathcal M:=(i(\omega I+\Omega\Ncal)-\mathcal R)^{-1}$, so
$\mathcal G(\omega)=\mathcal P\mathcal M\mathcal P^{-1}$ and hence
$\mathcal M=\mathcal P^{-1}\mathcal G(\omega)\mathcal P$, giving
$\norm{\mathcal M}\le\kappa(\mathcal P)\norm{\mathcal G(\omega)}$, i.e.
$\norm{\mathcal G(\omega)}\ge\norm{\mathcal M}/\kappa(\mathcal P)$. For any
matrix the resolvent norm is at least the reciprocal of the distance to the
spectrum, so
$\norm{(i(\omega+k\Omega)I-R)^{-1}}_2\ge 1/\min_j|i(\omega+k\Omega)-\lambda_j|$,
whence $\norm{\mathcal M}=\sup_k\norm{(i(\omega+k\Omega)I-R)^{-1}}_2\ge
1/\min_{j,k}|i\omega-\lambda_j+ik\Omega|=1/\gamma(\omega)$. Therefore
$\norm{\mathcal G(\omega)}\ge 1/(\gamma(\omega)\kappa(\mathcal P))\ge
1/(\gamma(\omega)\kappa_\star)$ (the last step since $\kappa(V)\ge1$). Finally, the
counterexample regime is real: with large $|\mu_j|$ one has $\delta(\omega)$
large yet $\gamma(\omega)=O(1)$, so $1/\delta(\omega)$ cannot upper-bound
$\norm{\mathcal G(\omega)}=\Theta(1/\gamma(\omega))$.
\end{proof}

\begin{proposition}[Resolvent as a periodic Green operator]\label{prop:green}
Under Assumption~\ref{ass:hyp} the bounded inverse
$\mathcal G(\omega):=\big(i(\omega I+\Omega\Ncal)-\Acal\big)^{-1}$ exists, and
for an EMP input $u(t)=\sum_\ell u_\ell e^{i(\omega+\ell\Omega)t}$ the state
$\mathbf x=\mathcal G(\omega)\Bcal\mathbf u$ is the harmonic-coefficient
sequence of the \emph{unique} $T$-quasiperiodic-of-rate-$\omega$ solution of
\eqref{eq:sys}, namely $x(t)=e^{i\omega t}\xi(t)$ where $\xi$ is the unique
$T$-periodic solution of
\begin{equation}\label{eq:bvp}
\dot\xi(t)=\big(A(t)-i\omega I\big)\xi(t)+\tilde b(t),\quad \xi(0)=\xi(T),\qquad
\tilde b(t):=\sum_\ell B(t)u_\ell e^{i\ell\Omega t}\,e^{-i\cdot 0}.
\end{equation}
\end{proposition}

\begin{proof}
Existence and boundedness of $\mathcal G(\omega)$ is Lemma~\ref{lem:spectrum}.
Write $x(t)=e^{i\omega t}\xi(t)$ with $\xi(t)=\sum_k x_k e^{ik\Omega t}$ (so
$\xi$ is $T$-periodic) and $u(t)=e^{i\omega t}\beta(t)$ with
$\beta(t)=\sum_\ell u_\ell e^{i\ell\Omega t}$ also $T$-periodic. Substituting
into $\dot x=A x+Bu$ gives
$e^{i\omega t}(\dot\xi+i\omega\xi)=e^{i\omega t}(A\xi+B\beta)$, i.e.
$\dot\xi=(A-i\omega I)\xi+B\beta$, which is \eqref{eq:bvp} with
$\tilde b=B\beta$. A $T$-periodic solution of this ODE is unique iff the
homogeneous problem $\dot\xi=(A-i\omega I)\xi$, $\xi(0)=\xi(T)$, has only the
trivial periodic solution. The transition matrix of $\dot\xi=(A-i\omega I)\xi$
is $e^{-i\omega(t-s)}\Phi_A(t,s)$, so its monodromy is $e^{-i\omega T}\Psi$, and
a nonzero periodic solution exists iff $1\in\spec(e^{-i\omega T}\Psi)$, i.e.
$e^{i\omega T}\in\spec(\Psi)$, which is excluded by \eqref{eq:hyp}. Hence
\eqref{eq:bvp} has a unique periodic solution; matching harmonics shows its
coefficient vector equals $\mathcal G(\omega)\Bcal\mathbf u$.
\end{proof}

Proposition~\ref{prop:green} is the structural pivot of the paper: \emph{the
infinite resolvent is the harmonic transform of a finite-dimensional periodic
two-point boundary-value problem.} This is what makes the monodromy algorithm
(\S\ref{sec:mono}) possible and exact.

\section{Algorithm I: structured Toeplitz truncation with a sharp error bound}
\label{sec:trunc}

The most direct method truncates $\Z$ to a window
$\mathbb{I}_M=\{-M,\dots,M\}$ and inverts the resulting finite matrix. The
contribution of this section is a \emph{provable, computable} error bound and an
exploitation of the band/decay structure.

For a window radius $M$ let $\Acal_M,\Bcal_M,\Ccal_M,\Dcal_M$ be the central
$(2M{+}1)$-block sections of \eqref{eq:Toep}, and
$\Ncal_M=\diag(kI_n)_{|k|\le M}$. Define the truncated HTF block
\begin{equation}\label{eq:HTFtrunc}
\Hcal^{(M)}(\omega)_{k\ell}
=\big[\Ccal_M\big(i(\omega I+\Omega\Ncal_M)-\Acal_M\big)^{-1}\Bcal_M+\Dcal_M\big]_{k\ell},
\quad |k|,|\ell|\le M_0,
\end{equation}
where $M_0<M$ is the requested number of \emph{output} harmonics.

\begin{algorithm}[H]
\caption{Truncated harmonic-balance HTF at one frequency}\label{alg:trunc}
\begin{algorithmic}[1]
\Require Fourier blocks $\{\hat A_k,\hat B_k,\hat C_k,\hat D_k\}$, frequency $\omega$,
output radius $M_0$, window radius $M\ge M_0$.
\Ensure Blocks $\Hcal^{(M)}(\omega)_{k\ell}$, $|k|,|\ell|\le M_0$.
\State Assemble block-Toeplitz $\Acal_M,\Bcal_M,\Ccal_M,\Dcal_M$ and $\Ncal_M$.
\State Form $W\gets i(\omega I+\Omega\Ncal_M)-\Acal_M$ (size $(2M{+}1)n$).
\State Factor $W=LU$ once (banded if data band-limited: bandwidth $(2K{+}1)n$).
\State Solve $WX=\Bcal_M$ for $X$ (the central columns suffice if $M_0<M$).
\State \Return $(\Ccal_M X+\Dcal_M)_{k\ell}$ for $|k|,|\ell|\le M_0$.
\end{algorithmic}
\end{algorithm}

\subsection{Numerical reliability}
$W=i(\omega I+\Omega\Ncal_M)-\Acal_M$ is a finite matrix; its conditioning is
governed by the generator gap $\gamma(\omega)$ in \eqref{eq:gengap}
(Proposition~\ref{prop:resnorm}: $\norm{W^{-1}}\le\kappa_\star/\gamma(\omega)$ up to the finite-section correction), \emph{not} by $\delta(\omega)$. Step~3 uses partial pivoting
(backward-stable). We never form the doubly-infinite operator, and for
band-limited data $W$ is banded so the factorization is stable and cheap.

\subsection{Error bound}
We prove a clean exponential bound on the truncation error. The decisive
structural fact is that the resolvent of $i\Omega\Ncal-\Acal$ has
exponentially decaying off-diagonal blocks; this is a quantitative form of the
Combes--Thomas / Paley--Wiener decay estimate.

\begin{lemma}[Off-diagonal decay of the resolvent]\label{lem:decay}
Assume the data are \emph{band-limited}, $\hat A_k=0$ for $|k|>K$ (the general
case is treated in Remark~\ref{rem:general}). Let
$G(\omega)=\mathcal G(\omega)$ with blocks $G_{k\ell}\in\C^{n\times n}$. There
exist constants $C_0>0$ and a rate $\alpha>0$, both computable from the data
and the \emph{generator gap} $\gamma(\omega)$ of \eqref{eq:gengap}, such that
\begin{equation}\label{eq:decay}
\norm{G_{k\ell}}_2\le C_0\,e^{-\alpha|k-\ell|},\qquad k,\ell\in\Z .
\end{equation}
Explicitly one may take any
\begin{equation}\label{eq:alpharange}
\alpha\in\Big(0,\ \tfrac1K\operatorname{arcsinh}\!\big(\gamma(\omega)/(4K\,\kappa_\star\,\norm{A}_{L^\infty})\big)\Big),
\end{equation}
where $\kappa_\star=\kappa(\mathcal P)\kappa(V)$ is the resolvent constant of
Proposition~\ref{prop:resnorm}, and a corresponding finite $C_0=C_0(\alpha,\omega)$.
The rate is governed by $\gamma(\omega)$, \emph{not} by $\delta(\omega)$.
\end{lemma}

\begin{proof}
For band-limited data $\Acal$ is block-banded with half-bandwidth $K$:
$(\Acal)_{k\ell}=0$ for $|k-\ell|>K$. Fix $\alpha>0$ and let
$E_\alpha=\diag(e^{\alpha k}I_n)_{k}$ (an unbounded but invertible weight). The
conjugated resolvent is
$E_\alpha^{-1}G E_\alpha=\big(i(\omega I+\Omega\Ncal)-E_\alpha^{-1}\Acal E_\alpha\big)^{-1}$,
because $\Ncal$ commutes with $E_\alpha$. The weighted operator
$\Acal_\alpha:=E_\alpha^{-1}\Acal E_\alpha$ has blocks
$(\Acal_\alpha)_{k\ell}=e^{-\alpha(k-\ell)}\hat A_{k-\ell}$, still supported in
$|k-\ell|\le K$, hence bounded with
$\norm{\Acal_\alpha-\Acal}\le \sum_{1\le|j|\le K}(e^{\alpha|j|}-1)\norm{\hat A_j}_2
\le (e^{\alpha K}-1)\,2K\max_j\norm{\hat A_j}_2$, which by
$\max_j\norm{\hat A_j}_2\le\norm{A}_{L^\infty}$ and $e^{\alpha K}-1\le
2\sinh(\alpha K)$ (valid for $\alpha K\le\log 2$, the regime of interest) is
$\le 4K\norm{A}_{L^\infty}\sinh(\alpha K)$. The unperturbed
resolvent is $G(\omega)=(i(\omega I+\Omega\Ncal)-\Acal)^{-1}$, whose norm is
bounded \emph{by the generator gap} via Proposition~\ref{prop:resnorm}:
\[
\norm{G}\le \kappa_\star/\gamma(\omega).
\]
(This is the correction of the earlier draft, which erroneously used
$\norm{G}\le c_\delta/\delta(\omega)$; that bound is false when the Floquet
multipliers are large, see the Warning after \eqref{eq:gap} and
Proposition~\ref{prop:resnorm}.) By the
Neumann/perturbation bound, if
$\norm{G}\,\norm{\Acal_\alpha-\Acal}<1$, i.e.
\begin{equation}\label{eq:alphacond}
4K\norm{A}_{L^\infty}\sinh(\alpha K)<\gamma(\omega)/\kappa_\star,
\end{equation}
then $E_\alpha^{-1}GE_\alpha=(I-G(\Acal-\Acal_\alpha))^{-1}G$ is bounded with norm
$\le \norm{G}/(1-\norm{G}\norm{\Acal-\Acal_\alpha})=:C_0$. Reading the
$(k,\ell)$ block gives
$\norm{(E_\alpha^{-1}GE_\alpha)_{k\ell}}_2=e^{-\alpha(k-\ell)}\norm{G_{k\ell}}_2$,
hence $\norm{G_{k\ell}}_2=e^{\alpha(k-\ell)}\norm{(E_\alpha^{-1}GE_\alpha)_{k\ell}}_2
\le C_0 e^{\alpha(k-\ell)}$; running the identical argument with the weight
$E_{-\alpha}$ gives $\norm{G_{k\ell}}_2\le C_0 e^{-\alpha(k-\ell)}$, and
combining the two yields $\norm{G_{k\ell}}_2\le C_0 e^{-\alpha|k-\ell|}$, which is
\eqref{eq:decay}. The admissible range of $\alpha$ is exactly the solvability of
\eqref{eq:alphacond}, i.e.
$\alpha<\tfrac1K\operatorname{arcsinh}\!\big(\gamma(\omega)/(4K\kappa_\star\norm{A}_{L^\infty})\big)$,
which is \eqref{eq:alpharange}.
\end{proof}

\begin{lemma}[Uniform finite-section stability and decay]\label{lem:finsec}
Assume band-limited data ($\hat A_k=0$, $|k|>K$) and Assumption~\ref{ass:hyp}.
Then there exist $M_\star\in\N$ and a constant $c_\star<\infty$, both computable
from the data and $\gamma(\omega)$, such that for every window radius $M\ge
M_\star$ the central section $L_M=i(\omega I+\Omega\Ncal_M)-\Acal_M$ is invertible
on $\operatorname{ran}\Pi_M$ with
\begin{equation}\label{eq:unifstab}
\norm{L_M^{-1}}\;\le\; c_\star,\qquad\text{uniformly in }M\ge M_\star .
\end{equation}
Moreover the inverses $G^{(M)}=L_M^{-1}\Pi_M$ satisfy the off-diagonal decay
\begin{equation}\label{eq:decayFS}
\norm{(G^{(M)})_{k\ell}}_2\le C_0'\,e^{-\alpha'|k-\ell|}\qquad(|k|,|\ell|\le M),
\end{equation}
with constants $C_0',\alpha'>0$ \emph{independent of $M$}, for any
$\alpha'\in\big(0,\tfrac1K\operatorname{arcsinh}(1/(4Kc_\star\norm{A}_{L^\infty}))\big)$.
\end{lemma}

\begin{proof}
\emph{Step 1 (tail diagonal dominance).} The diagonal blocks of $L$ are
$D_k:=i(\omega+k\Omega)I-\hat A_0$ and the off-diagonal mass in any block-row is
$\sum_{1\le|j|\le K}\norm{\hat A_j}_2\le 2K\norm{A}_{L^\infty}=:a$. For
$|k|\ge k_0:=\big\lceil(\,\|\hat A_0\|_2+2a\,)/\Omega+|\omega|/\Omega\big\rceil$ we
have $\norm{D_k^{-1}}_2^{-1}=\min_j|i(\omega+k\Omega)-\mathrm{spec}(\hat A_0)_j|
\ge |k|\Omega-|\omega|-\norm{\hat A_0}_2\ge 2a$, so the block-row is strictly
diagonally dominant with dominance factor $\ge 2$:
$\norm{D_k^{-1}}_2\cdot a\le 1/2$.

\emph{Step 2 (collar Schur complement).} Fix any $M\ge M_\star:=k_0+K$. Split the
window indices $\mathbb I_M$ into the \emph{core} $\mathcal I=\{|k|\le k_0\}$ and
the two \emph{tails} $\mathcal T=\{k_0<|k|\le M\}$. Order $L_M=\left(\begin{smallmatrix}
L_{\mathcal I\mathcal I}&L_{\mathcal I\mathcal T}\\ L_{\mathcal T\mathcal I}&
L_{\mathcal T\mathcal T}\end{smallmatrix}\right)$. The tail block $L_{\mathcal T\mathcal T}$
is block-tridiagonal-banded (bandwidth $K$) with diagonally dominant rows by
Step~1, hence invertible with $\norm{L_{\mathcal T\mathcal T}^{-1}}\le 1/a$
(standard bound for block-diagonally-dominant matrices: if
$\norm{D_k^{-1}}\,(\text{row off-diag mass})\le\theta<1$ for all rows, then the
matrix is invertible and $\norm{(\cdot)^{-1}}\le \sup_k\norm{D_k^{-1}}/(1-\theta)
\le (1/2a)/(1/2)=1/a$). This bound is \emph{uniform in $M$} because it only uses
the per-row dominance, which holds for every tail row regardless of $M$.

The Schur complement of $L_{\mathcal T\mathcal T}$ in $L_M$ is
$S_M:=L_{\mathcal I\mathcal I}-L_{\mathcal I\mathcal T}L_{\mathcal T\mathcal T}^{-1}
L_{\mathcal T\mathcal I}$, a fixed-size ($(2k_0{+}1)n$) matrix. We claim
$S_M\to S_\infty$ as $M\to\infty$, where $S_\infty$ is the Schur complement of
the (invertible, by Step~1) infinite tail operator $L_{\mathcal T_\infty
\mathcal T_\infty}$ ($\mathcal T_\infty=\{|k|>k_0\}$) in $L$; indeed
$L_{\mathcal I\mathcal T}$ has only the $O(K)$ columns adjacent to the core
nonzero (bandedness), $L_{\mathcal T\mathcal I}$ likewise, and
$L_{\mathcal T\mathcal T}^{-1}$ converges entrywise to
$L_{\mathcal T_\infty\mathcal T_\infty}^{-1}$ on those finitely many relevant
entries with geometric rate (again diagonal dominance gives geometric decay of
the banded inverse away from the diagonal, so the finite-tail inverse and the
infinite-tail inverse agree up to $O(e^{-c(M-k_0)})$ on the core-adjacent strip).
Now $S_\infty$ is invertible: $L$ is invertible (Lemma~\ref{lem:spectrum}), and
by the Schur-complement/Banachiewicz identity for the $2\times2$ block operator
$L$ with invertible tail block, $L^{-1}$ exists iff $S_\infty$ is invertible;
hence $S_\infty^{-1}$ exists with $\norm{S_\infty^{-1}}=\norm{\Pi_{\mathcal I}L^{-1}
\Pi_{\mathcal I}\text{-block-inverse}}<\infty$ (explicitly,
$S_\infty^{-1}=\big(\Pi_{\mathcal I}L^{-1}\Pi_{\mathcal I}\big)$ restricted to
the core, which is bounded by $\norm{L^{-1}}\le\kappa_\star/\gamma(\omega)$).
By convergence $S_M\to S_\infty$, for $M$ large enough
$\norm{S_M^{-1}}\le 2\norm{S_\infty^{-1}}\le 2\kappa_\star/\gamma(\omega)=:c_S$.

\emph{Step 3 (block inversion bound).} With both $L_{\mathcal T\mathcal T}$ and
$S_M$ invertible, the standard $2\times2$ block-inverse formula expresses
$L_M^{-1}$ in terms of $L_{\mathcal T\mathcal T}^{-1}$ (norm $\le 1/a$, uniform)
and $S_M^{-1}$ (norm $\le c_S$, uniform) and the bounded coupling blocks
($\norm{L_{\mathcal I\mathcal T}},\norm{L_{\mathcal T\mathcal I}}\le a$). Hence
$\norm{L_M^{-1}}\le c_\star$ with an explicit
$c_\star=c_\star(a,c_S,1/a)$ \emph{independent of $M$}, proving
\eqref{eq:unifstab}.

\emph{Step 4 (uniform decay).} With the uniform bound \eqref{eq:unifstab} in
hand, repeat the weighted-conjugation/perturbation argument of
Lemma~\ref{lem:decay} \emph{on the finite matrix} $L_M$: the diagonal weight
$E_{\pm\alpha'}$ restricted to $\mathbb I_M$ commutes with $\Ncal_M$, the banded
$\Acal_M$ satisfies $\norm{E_{\alpha'}^{-1}\Acal_M E_{\alpha'}-\Acal_M}\le
4K\norm{A}_{L^\infty}\sinh(\alpha' K)$ \emph{exactly as in the infinite case}
(the weight is diagonal and the band structure is identical), and the
unperturbed inverse bound is now \eqref{eq:unifstab}. The Neumann condition
$c_\star\cdot 4K\norm{A}_{L^\infty}\sinh(\alpha' K)<1$ defines the stated
$\alpha'$-range and yields
$\norm{(G^{(M)})_{k\ell}}_2\le C_0'e^{-\alpha'|k-\ell|}$ with
$C_0'=c_\star/(1-c_\star\,4K\norm{A}_{L^\infty}\sinh(\alpha'K))$, both
$M$-independent. This is \eqref{eq:decayFS}.
\end{proof}

\begin{theorem}[Sharp truncation error bound]\label{thm:trunc}
Under the hypotheses of Lemmas~\ref{lem:decay} and \ref{lem:finsec}, for
$M\ge\max(M_0,M_\star)$ and $|k|,|\ell|\le M_0$,
\begin{equation}\label{eq:trerr}
\norm{\Hcal(\omega)_{k\ell}-\Hcal^{(M)}(\omega)_{k\ell}}_2
\;\le\; C_1\,e^{-\bar\alpha\,(M-M_0)},
\end{equation}
where $\bar\alpha:=\tfrac12\min(\alpha_{\max},\alpha'_{\max})$ is the fixed
(interior, hence with finite constants) rate, $\alpha_{\max},\alpha'_{\max}$ being
the endpoints of the admissible ranges in Lemmas~\ref{lem:decay},\ref{lem:finsec},
and $C_1$ is computable and \emph{finite} (because $\bar\alpha$ is bounded away
from the endpoints where $C_0,C_0'$ blow up) from
$\bar C=\max(C_0,C_0')$, $\norm{B}_{L^\infty}$, $\norm{C}_{L^\infty}$,
$\bar\alpha$ and $K$. In particular, to obtain absolute tolerance $\varepsilon$
on the output harmonics $|k|\le M_0$ it suffices to take
\begin{equation}\label{eq:Mchoice}
M\;=\;\max(M_0,M_\star)+\Big\lceil \tfrac1{\bar\alpha}\log\tfrac{C_1}{\varepsilon}\Big\rceil .
\end{equation}
\end{theorem}

\begin{proof}
Both $\Hcal$ and $\Hcal^{(M)}$ have the same $\Dcal$ term, which cancels, and
the same $\Ccal,\Bcal$ blocks for the relevant indices. Write the exact resolvent
$G=\mathcal G(\omega)$ and the truncated resolvent
$G^{(M)}=(i(\omega I+\Omega\Ncal_M)-\Acal_M)^{-1}$ (extended by $0$ outside the
window). It suffices to bound $G_{k\ell}-G^{(M)}_{k\ell}$ for $|k|,|\ell|\le M_0$
and then multiply by the bounded $\Ccal,\Bcal$ Toeplitz tails.

Let $\Pi_M$ be the orthogonal projection of $\ell^2(\Z;\C^n)$ onto the window
$\mathbb I_M=\{-M,\dots,M\}$ and $\Pi_M^\perp=I-\Pi_M$. Write
$L:=i(\omega I+\Omega\Ncal)-\Acal$, so $G=L^{-1}$, and let
$L_M:=\Pi_M L\,\Pi_M$ be the central section (its matrix is
$i(\omega I+\Omega\Ncal_M)-\Acal_M=W$ from Algorithm~\ref{alg:trunc}). By
Lemma~\ref{lem:finsec}, for all $M\ge M_\star$ the section $L_M$ is invertible on
$\operatorname{ran}\Pi_M$ with the \emph{uniform} bound $\norm{L_M^{-1}}\le
c_\star$ \eqref{eq:unifstab}; this is the genuine finite-section stability, proved
there by a tail diagonal-dominance/Schur-complement argument rather than assumed.
Define
$G^{(M)}:=L_M^{-1}\Pi_M$ (extended by $0$ off the window); this is the operator
inverted in \eqref{eq:HTFtrunc}.

We use the finite-section resolvent identity. On $\operatorname{ran}\Pi_M$,
$L_M = \Pi_M L \Pi_M = \Pi_M L - \Pi_M L \Pi_M^\perp$, hence
$\Pi_M L = L_M + \Pi_M L \Pi_M^\perp$. Right-multiplying by $G=L^{-1}$ and
left-multiplying by $L_M^{-1}$ gives, on $\operatorname{ran}\Pi_M$,
\[
L_M^{-1}\Pi_M = G + L_M^{-1}\,\Pi_M L\,\Pi_M^\perp\,G ,
\qquad\text{i.e.}\qquad
G^{(M)}-\Pi_M G\Pi_M
= L_M^{-1}\,\big(\Pi_M L\,\Pi_M^\perp\big)\,G\,\Pi_M .
\]
The operator $\Pi_M L\,\Pi_M^\perp$ equals $-\Pi_M\Acal\Pi_M^\perp$ (the
diagonal part $i(\omega I+\Omega\Ncal)$ commutes with the projections and so
contributes nothing to the off-diagonal block $\Pi_M(\cdot)\Pi_M^\perp$).
Because $\Acal$ is block-banded with half-bandwidth $K$, the block
$\Pi_M\Acal\Pi_M^\perp$ has entries $(\Acal)_{rj}=\hat A_{r-j}$ supported on the
\emph{boundary strip}
$\mathcal S:=\{(r,j): |r|\le M<|j|,\ 1\le |r-j|\le K\}$, which contains $O(K^2)$
nonzero blocks (only $r\in\{M-K+1,\dots,M\}\cup\{-M,\dots,-M+K-1\}$ couple
across the window edge). Reading the $(k,\ell)$ block of the identity for
$|k|,|\ell|\le M_0$ (so $\Pi_M G\Pi_M$ contributes exactly $G_{k\ell}$ and
$G^{(M)}_{k\ell}$ is the truncated block), and using
$L_M^{-1}=G^{(M)}$ on the range,
\[
G^{(M)}_{k\ell}-G_{k\ell}
= -\sum_{(r,j)\in\mathcal S} G^{(M)}_{kr}\,\hat A_{r-j}\,G_{j\ell},
\]
whence
\[
\norm{(G-G^{(M)})_{k\ell}}_2
\;\le\; \sum_{(r,j)\in\mathcal S}
\norm{G^{(M)}_{kr}}_2\,\norm{\hat A_{r-j}}_2\,\norm{G_{j\ell}}_2 .
\]
Apply the off-diagonal decay bound \eqref{eq:decay} to the exact resolvent,
$\norm{G_{j\ell}}_2\le C_0 e^{-\alpha|j-\ell|}$, and the finite-section decay bound \eqref{eq:decayFS} of
Lemma~\ref{lem:finsec} to the truncated inverse,
$\norm{G^{(M)}_{kr}}_2\le C_0'e^{-\alpha'|k-r|}$, with constants $C_0',\alpha'$
that are \emph{independent of $M$} (this is the content of
Lemma~\ref{lem:finsec}, Step~4, which does \emph{not} transfer the infinite
conjugation by $\mathcal P$---that would fail since $\mathcal P$ does not commute
with $\Pi_M$---but instead re-runs the weighted perturbation argument directly on
the finite matrix using the uniform stability \eqref{eq:unifstab}). On the strip $\mathcal S$ we have
$|k|\le M_0$ and $r\in\{M-K+1,\dots,M\}\cup\{-M,\dots,-M+K-1\}$, hence
$|k-r|\ge M-K+1-M_0$; and $|\ell|\le M_0$ with $|j|>M$, hence
$|j-\ell|\ge M+1-M_0$. Set $\bar\alpha:=\min(\alpha,\alpha')$ and
$\bar C:=\max(C_0,C_0')$ (so both the exact and the finite-section resolvent obey
the bound $C_0,C_0'\le\bar C$ with rate $\ge\bar\alpha$). Therefore each summand is
$\le \bar C^2\,\norm{A}_{L^\infty}\,e^{-\bar\alpha(|k-r|+|j-\ell|)}
\le \bar C^2\,\norm{A}_{L^\infty}\,e^{-\bar\alpha(2(M-M_0)-K+2)}$. The
geometric sum over the $O(K^2)$ boundary couplings is then bounded by
$C_1' e^{-\bar\alpha(M-M_0)}$ with
$C_1'=K^2\norm{A}_{L^\infty}\bar C^2 e^{2\bar\alpha K}/(1-e^{-\bar\alpha})^2$
(absorbing the remaining $e^{-\bar\alpha(M-M_0)}\le 1$ factor and the geometric
constants). Finally
multiply by the bounded output/input maps,
$\norm{\Ccal}\le\norm{C}_{L^\infty}$, $\norm{\Bcal}\le\norm{B}_{L^\infty}$,
to obtain \eqref{eq:trerr} with
$C_1=\norm{C}_{L^\infty}\norm{B}_{L^\infty}C_1'$. Solving
$C_1 e^{-\bar\alpha(M-M_0)}\le\varepsilon$ gives \eqref{eq:Mchoice} (with $\alpha$ there read as $\bar\alpha$).
\end{proof}

\begin{proposition}[Constructive bound for non-band-limited data]\label{prop:general}
Suppose the data are not band-limited but have summable Fourier tails. Fix a
truncation level $K$ and let $A_K(t):=\sum_{|k|\le K}\hat A_k e^{ik\Omega t}$ be
the band-limited symbol, with consistency error
$\eta_K:=\norm{A-A_K}_{L^\infty}\le\sum_{|k|>K}\norm{\hat A_k}_2$
(computable; e.g. $\le c\,r^{K+1}/(1-r)$ for $\norm{\hat A_k}_2\le c r^{|k|}$, or
$\le c'K^{-s+1}$ for $\norm{\hat A_k}_2\le c|k|^{-s}$). Let $\Acal_K$ be the
corresponding (banded) Laurent operator and $G_K=(i(\omega I+\Omega\Ncal)-\Acal_K)^{-1}$.
If $\eta_K\,\norm{G}<1$ (which holds for $K$ large enough since
$\eta_K\to0$), then $G_K$ exists, $\norm{G_K}\le\norm{G}/(1-\eta_K\norm{G})$,
and
\begin{equation}\label{eq:Gconsist}
\norm{G-G_K}\le \frac{\norm{G}^2\eta_K}{1-\eta_K\norm{G}}
\le \frac{(\kappa_\star/\gamma(\omega))^2\,\eta_K}{1-\eta_K\kappa_\star/\gamma(\omega)},
\end{equation}
so $\norm{\Hcal(\omega)-\Hcal_K(\omega)}\le\norm{C}_{L^\infty}\norm{B}_{L^\infty}
\norm{G-G_K}$. Applying Lemmas~\ref{lem:decay}--\ref{lem:finsec} and
Theorem~\ref{thm:trunc} to the band-limited $A_K$ then gives a fully
\emph{constructive} two-parameter error bound
\begin{equation}\label{eq:twoparam}
\norm{\Hcal(\omega)_{k\ell}-\Hcal_K^{(M)}(\omega)_{k\ell}}_2
\le \underbrace{C_1\,e^{-\bar\alpha_K(M-M_0)}}_{\text{truncation in }M}
+\underbrace{\norm{C}_{L^\infty}\norm{B}_{L^\infty}\,\tfrac{\norm{G}^2\eta_K}{1-\eta_K\norm{G}}}_{\text{symbol truncation in }K},
\end{equation}
where $\bar\alpha_K$ is the rate of Theorem~\ref{thm:trunc} for the symbol $A_K$.
Both terms are computable; choosing $K$ to make the second $\le\varepsilon/2$ and
$M$ to make the first $\le\varepsilon/2$ achieves total error $\le\varepsilon$.
\end{proposition}

\begin{proof}
Write $\Acal=\Acal_K+\Delta$ with $\norm\Delta=\norm{A-A_K}_{L^\infty}=\eta_K$
(Lemma~\ref{lem:bounded}). Then
$i(\omega I+\Omega\Ncal)-\Acal_K=(i(\omega I+\Omega\Ncal)-\Acal)+\Delta
=G^{-1}(I+G\Delta)$, which is invertible whenever $\norm{G\Delta}\le\eta_K\norm G<1$,
with $G_K=(I+G\Delta)^{-1}G$ and the Neumann bounds
$\norm{G_K}\le\norm G/(1-\eta_K\norm G)$ and
$\norm{G-G_K}=\norm{G_K\Delta G}\le\norm{G_K}\eta_K\norm G$, giving
\eqref{eq:Gconsist} after inserting $\norm G\le\kappa_\star/\gamma(\omega)$
(Proposition~\ref{prop:resnorm}). The composite bound \eqref{eq:twoparam} is the
triangle inequality
$\norm{\Hcal-\Hcal_K^{(M)}}\le\norm{\Hcal-\Hcal_K}+\norm{\Hcal_K-\Hcal_K^{(M)}}$,
the first term being the symbol-consistency error just bounded and the second the
band-limited truncation error of Theorem~\ref{thm:trunc} applied to $A_K$.
\end{proof}

\begin{remark}[Qualitative inverse-closedness]\label{rem:general}
The constructive route of Proposition~\ref{prop:general} is the recommended one
because it yields a computable $M,K$. For completeness we note the qualitative
companion fact: the class of operators with
$\norm{(\Acal)_{k\ell}}_2\le c\,\rho^{|k-\ell|}$ ($0<\rho<1$) is a Banach algebra
(a Jaffard/Baskakov–Gohberg–Sjöstrand class) closed under inversion, so if
$\Acal$ has geometric off-diagonal decay and $G$ is bounded then
$\norm{G_{k\ell}}_2\le C_0\tilde\rho^{|k-\ell|}$ for some $\tilde\rho\in(\rho,1)$
(Bochner–Phillips/Jaffard inverse-closedness). This gives the qualitative decay
$\norm{\Hcal_{k\ell}-\Hcal^{(M)}_{k\ell}}_2\le C_1\tilde\rho^{M-M_0}$ but with a
non-explicit $\tilde\rho$; Proposition~\ref{prop:general} replaces it by explicit
constants at the cost of the extra symbol-truncation parameter $K$.
\end{remark}

\section{Algorithm II: monodromy/Floquet evaluation (recommended)}
\label{sec:mono}

We now give the method that \emph{exactly} avoids the infinite operator. By
Proposition~\ref{prop:green}, computing the $\ell$-th block-column of
$\Hcal(\omega)$ (the response to a unit input in harmonic $\ell$ of input
channel) is solving the finite-dimensional periodic BVP \eqref{eq:bvp} with
forcing $\tilde b(t)=B(t)e^{i\ell\Omega t}e_c$ ($e_c$ a unit input vector),
extracting the periodic state $\xi(t)$, forming $C(t)\xi(t)+D(t)e^{i\ell\Omega t}e_c$,
and reading its Fourier coefficients $y_k$. Crucially this is independent of any
harmonic truncation of $A$: $\xi$ is computed by integrating the genuine
$n\times n$ ODE.

\begin{lemma}[Closed-form periodic BVP solution via monodromy]\label{lem:monoform}
Let $\Phi(t):=e^{-i\omega t}\Phi_A(t,0)$ be the transition matrix of
$\dot\xi=(A(t)-i\omega I)\xi$, so $\Phi(0)=I$ and $\Phi(T)=e^{-i\omega T}\Psi$.
Under \eqref{eq:hyp} the unique $T$-periodic solution of \eqref{eq:bvp} is
\begin{equation}\label{eq:xiform}
\xi(t)=\Phi(t)\Big[(I-\Phi(T))^{-1}\Phi(T)\,v(T)+v(t)\Big],\qquad
v(t):=\int_0^t\Phi(s)^{-1}\tilde b(s)\,ds .
\end{equation}
\end{lemma}

\begin{proof}
By variation of constants, every solution of $\dot\xi=(A-i\omega I)\xi+\tilde b$
with $\xi(0)=\xi_0$ is $\xi(t)=\Phi(t)(\xi_0+v(t))$, since
$\dot\xi=\dot\Phi(\xi_0+v)+\Phi\dot v
=(A-i\omega I)\Phi(\xi_0+v)+\Phi\Phi^{-1}\tilde b=(A-i\omega I)\xi+\tilde b$, and
$\xi(0)=\Phi(0)\xi_0=\xi_0$. Periodicity $\xi(T)=\xi(0)$ reads
$\Phi(T)(\xi_0+v(T))=\xi_0$, i.e. $(I-\Phi(T))\xi_0=\Phi(T)v(T)$. By
\eqref{eq:hyp}, $1\notin\spec(\Phi(T))=e^{-i\omega T}\spec(\Psi)$, so $I-\Phi(T)$
is invertible and $\xi_0=(I-\Phi(T))^{-1}\Phi(T)v(T)$. Substituting gives
\eqref{eq:xiform}; uniqueness follows from invertibility of $I-\Phi(T)$.
\end{proof}

Formula \eqref{eq:xiform} was verified numerically to reproduce the HTF blocks
to machine precision on random LTP systems. The full algorithm follows.

\begin{algorithm}[H]
\caption{Monodromy/Floquet HTF at a single frequency $\omega$}\label{alg:mono}
\begin{algorithmic}[1]
\Require Data $A(\cdot),B(\cdot),C(\cdot),D(\cdot)$ ($T$-periodic), frequency $\omega$
with $e^{i\omega T}\notin\spec(\Psi)$, output radius $M_0$, integration tol $\eta$.
\Ensure Blocks $\Hcal(\omega)_{k\ell}$, $|k|,|\ell|\le M_0$.
\State Integrate the $n\times n$ matrix ODE $\dot\Phi=(A(t)-i\omega I)\Phi$,
$\Phi(0)=I$, on $[0,T]$ with a stiffly-stable, error-controlled solver (tol
$\eta$); store a dense interpolant $\Phi(t)$. Set $\Phi_T:=\Phi(T)$.
\For{each input harmonic/channel index $\ell$ (one unit forcing each)}
  \State Compute $v(T)=\int_0^T\Phi(s)^{-1}B(s)e^{i\ell\Omega s}\,ds\;e_c$ and the
  interpolant $v(t)$ by augmenting the ODE quadrature.
  \State Solve $(I-\Phi_T)\,\xi_0=\Phi_T v(T)$ (an $n\times n$ linear solve).
  \State Form $\xi(t)=\Phi(t)(\xi_0+v(t))$ and
  $w(t)=C(t)\xi(t)+D(t)e^{i\ell\Omega t}e_c$ on a grid of $N_F$ points,
  $N_F$ chosen per Theorem~\ref{thm:mono} to make aliasing $\le\varepsilon$.
  \State $\Hcal(\omega)_{k\ell}\gets \mathrm{FFT}_k\{w\}$ for $|k|\le M_0$
  (the $k$-th Fourier coefficient of $w$).
\EndFor
\State \Return $\{\Hcal(\omega)_{k\ell}\}_{|k|,|\ell|\le M_0}$.
\end{algorithmic}
\end{algorithm}

\begin{theorem}[Correctness and error of Algorithm~\ref{alg:mono}]\label{thm:mono}
Suppose the integrator in Step~1 produces $\Phi$ with global error
$\norm{\Phi-\hat\Phi}_{L^\infty[0,T]}\le\eta$ and the quadrature in Step~3 has
error $\le\eta$. Note that the periodic state $\xi$ is in general \emph{not}
band-limited even when $A,B$ are: by Lemma~\ref{lem:decay} applied to the
coefficient sequence of $\xi$ (which equals a column of $G(\omega)\Bcal$, hence
inherits geometric off-diagonal decay), the Fourier coefficients $\hat\xi_k$
obey $\norm{\hat\xi_k}_2\le C_\xi e^{-\alpha|k|}$ for the rate $\alpha$ of
\eqref{eq:decay}; therefore $w(t)=C(t)\xi(t)+D(t)e^{i\ell\Omega t}e_c$ has
geometrically decaying harmonics $\norm{\hat w_k}_2\le C_w e^{-\alpha|k|}$ for
$|k|$ large. The DFT on $N_F$ equispaced nodes then aliases harmonics of order
$\ge N_F-M_0$ into $|k|\le M_0$. Then the
computed blocks satisfy, for $|k|,|\ell|\le M_0$,
\begin{equation}\label{eq:monoerr}
\norm{\Hcal(\omega)_{k\ell}-\widehat{\Hcal}(\omega)_{k\ell}}_2
\le \kappa_\Phi(\omega)\,\big(\norm{C}_{L^\infty}\norm{\xi}_{L^\infty}+\text{quad}\big)\,\eta
\;+\;\text{(aliasing)},
\end{equation}
where $\kappa_\Phi(\omega)=O\!\big(\norm{(I-\Phi_T)^{-1}}\big)
=O(\kappa(V)\,\delta(\omega)^{-1})$ is the condition number of the BVP linear solve in
Step~4, and the aliasing term is bounded by
$\sum_{|j|\ge N_F-M_0}\norm{\hat w_j}_2\le C_w\,e^{-\alpha(N_F-M_0)}/(1-e^{-\alpha})$,
which decays geometrically in $N_F$; choosing
$N_F=M_0+\lceil\tfrac1\alpha\log\tfrac{C_w}{\varepsilon(1-e^{-\alpha})}\rceil$
makes it $\le\varepsilon$. There is \emph{no} harmonic-truncation error in $A$:
the dependence on $A$ is through the exact ODE.
\end{theorem}

\begin{proof}
Exactness of the analytic step is Lemma~\ref{lem:monoform} together with
Proposition~\ref{prop:green}, which identifies the Fourier coefficients of
$w(t)=C(t)\xi(t)+D(t)e^{i\ell\Omega t}e_c$ with the blocks
$\Hcal(\omega)_{k\ell}$. We bound the three numerical errors. (i) \emph{ODE
error}: by Gr\"onwall and the variation-of-constants formula, a perturbation
$\hat\Phi=\Phi+\Delta$, $\norm\Delta\le\eta$, propagates to
$\hat\xi=\xi+O(\eta)$ with constant
$\le \norm{(I-\Phi_T)^{-1}}(1+\norm{\Phi}_\infty^2)$, hence the linear-solve
condition number $\kappa_\Phi(\omega)$ enters as stated; by \eqref{eq:hyp},
$\norm{(I-\Phi_T)^{-1}}\le \kappa(V)/\dist(1,\spec(\Phi_T))=
\kappa(V)/\delta(\omega)$, where $\Psi=V\Lambda V^{-1}$ and
$\kappa(V)=\norm V\norm{V^{-1}}$ is the eigenvector conditioning of $\Psi$
(for non-normal $\Psi$ the factor $\kappa(V)$ is genuinely necessary: the bound
$\norm{(I-\Phi_T)^{-1}}\le 1/\dist(1,\spec(\Phi_T))$ holds only for normal
$\Phi_T$). Here $\delta(\omega)$ \emph{is} the correct gap because this is a
\emph{distance-to-$1$} (BVP-solvability) quantity, distinct from the
lifted-resolvent norm of Algorithm~I which is governed by $\gamma(\omega)$. (ii) \emph{Quadrature error}: enters linearly in $\eta$
through $v$, multiplied by $\norm{\Phi}_\infty$ and the same BVP solve. (iii)
\emph{Aliasing}: the trapezoidal/DFT rule on $N_F$ equispaced nodes
returns $\widehat w_k^{\mathrm{DFT}}=\sum_{r\in\Z}\hat w_{k+rN_F}$, so the
aliasing error at index $|k|\le M_0$ is $\sum_{r\ne0}\hat w_{k+rN_F}$, bounded by
$\sum_{|j|\ge N_F-M_0}\norm{\hat w_j}_2$. As established in the theorem statement,
$\norm{\hat w_j}_2\le C_w e^{-\alpha|j|}$ (geometric decay of $\hat\xi$ from
Lemma~\ref{lem:decay}, multiplied by the bounded Toeplitz operator $\Ccal$ and
added to the single $\Dcal$-harmonic), whence the tail is
$\le C_w e^{-\alpha(N_F-M_0)}/(1-e^{-\alpha})$. For data whose harmonics decay
only algebraically, $\norm{\hat A_k}_2\le c|k|^{-s}$, the same step gives an
algebraic tail $O((N_F-M_0)^{-s+1})$. Summing (i)--(iii) gives
\eqref{eq:monoerr}.
\end{proof}

\begin{remark}[Why this is numerically reliable]
Algorithm~\ref{alg:mono} replaces the inversion of a $(2M{+}1)n$ operator by:
one $n\times n$ matrix ODE solve (error-controlled, backward-stable), $m$
quadratures, $m$ $n\times n$ linear solves, and FFTs. The only ill-conditioning
is the genuine, unavoidable one near poles, quantified by $\delta(\omega)^{-1}$;
this is identical to the LTI situation where the transfer function blows up near
a pole. Stiff/large monodromy is handled by integrating $\dot\Phi=(A-i\omega
I)\Phi$ with a stiff solver, or by the scaled/Magnus exponential integrators
that compute $\Phi$ in a numerically stable factored form
$\Phi(T)=\prod_j e^{H_j}$ (the Floquet factorization), avoiding overflow.
\end{remark}

\begin{lemma}[Relative error of the factored monodromy integrator]\label{lem:stiff}
Partition $[0,T]$ into $N_s$ panels $0=t_0<\dots<t_{N_s}=T$ and let
$\hat\Phi(t_{i+1},t_i)$ approximate the panel transition matrix
$\Phi(t_{i+1},t_i)$ with \emph{relative} error
$\norm{\hat\Phi(t_{i+1},t_i)-\Phi(t_{i+1},t_i)}\le\eta_{\mathrm{rel}}
\norm{\Phi(t_{i+1},t_i)}$ (achievable to $\eta_{\mathrm{rel}}=O(\varepsilon_{\mathrm{mach}})$
by a scaled-and-squared or Magnus exponential of the panel, since each panel
matrix has $O(1)$ norm when the panels resolve $A$). Composing,
$\hat\Phi(T,0)=\prod_i\hat\Phi(t_{i+1},t_i)$ satisfies the multiplicative bound
\begin{equation}\label{eq:stifferr}
\frac{\norm{\hat\Phi(T,0)-\Phi(T,0)}}{\prod_i\norm{\Phi(t_{i+1},t_i)}}
\le (1+\eta_{\mathrm{rel}})^{N_s}-1\le N_s\eta_{\mathrm{rel}}\,e^{N_s\eta_{\mathrm{rel}}}.
\end{equation}
Hence the \emph{relative} accuracy of $\Psi=\Phi(T,0)$ is $O(N_s\eta_{\mathrm{rel}})$
\emph{regardless of how large $\norm\Psi$ (or any partial product) is}: unstable
Floquet modes do not amplify the relative error, because the bound is normalized
by the product of panel norms rather than by a single Gr\"onwall exponential.
Consequently the absolute error $\norm{\hat\Phi-\Phi}$ scales with the true
$\norm\Phi$, so the downstream BVP solve $(I-\Phi_T)\xi_0=\Phi_T v(T)$ inherits
only the genuine conditioning $\kappa(V)/\delta(\omega)$ of
Theorem~\ref{thm:mono}, with no spurious exponential blow-up.
\end{lemma}

\begin{proof}
For a product of factors with $\hat F_i=F_i+E_i$, $\norm{E_i}\le\eta_{\mathrm{rel}}\norm{F_i}$,
a telescoping expansion gives
$\prod_i\hat F_i-\prod_i F_i=\sum_i(\prod_{j<i}\hat F_j)E_i(\prod_{j>i}F_j)$,
whence
$\norm{\prod\hat F_i-\prod F_i}\le\sum_i(\prod_{j<i}(1+\eta_{\mathrm{rel}})\norm{F_j})
\eta_{\mathrm{rel}}\norm{F_i}(\prod_{j>i}\norm{F_j})
\le\big((1+\eta_{\mathrm{rel}})^{N_s}-1\big)\prod_i\norm{F_i}$,
which is \eqref{eq:stifferr}; the final inequality is
$(1+x)^N-1\le Nx\,e^{Nx}$. Achievability of panel-wise relative accuracy by
scaling-and-squaring/Magnus is the standard backward-stability of the matrix
exponential. The BVP-solve conditioning statement is Theorem~\ref{thm:mono}.
\end{proof}

\section{Algorithm III: fast frequency sweep over a grid}\label{sec:lift}

A grid of $N_\omega$ frequencies is requested. Naively running
Algorithm~\ref{alg:mono} per frequency costs $N_\omega$ ODE solves. We reduce
this to \emph{one} ODE solve plus $O(N_\omega)$ cheap algebraic updates by
exploiting that the $\omega$-dependence of $\Phi(t)=e^{-i\omega t}\Phi_A(t,0)$ is
explicit.

\begin{proposition}[Frequency-shift factorization]\label{prop:shift}
Let $\Phi_A(t,0)$ be the $\omega$-independent transition matrix of
$\dot z=A(t)z$, computed once. Then for every $\omega$,
\begin{equation}\label{eq:shift}
\Phi(t)=e^{-i\omega t}\Phi_A(t,0),\quad
\Phi(T)=e^{-i\omega T}\Psi,\quad
\Phi(s)^{-1}\tilde b(s)=e^{i\omega s}\Phi_A(s,0)^{-1}B(s)e^{i\ell\Omega s}e_c .
\end{equation}
Consequently $v(T)=\int_0^T e^{i\omega s}\,g_\ell(s)\,ds$ with the
$\omega$-independent integrand $g_\ell(s)=\Phi_A(s,0)^{-1}B(s)e^{i\ell\Omega s}e_c$,
so $v(T)$ is the Fourier transform of $g_\ell$ evaluated at $-\omega$, and the
BVP solve uses only $\Psi$ and the scalar $e^{-i\omega T}$.
\end{proposition}

\begin{proof}
$\Phi$ solves $\dot\Phi=(A-i\omega I)\Phi$, $\Phi(0)=I$; the factor
$e^{-i\omega t}\Phi_A(t,0)$ solves it since
$\tfrac{d}{dt}(e^{-i\omega t}\Phi_A)=-i\omega e^{-i\omega t}\Phi_A+e^{-i\omega
t}A\Phi_A=(A-i\omega I)(e^{-i\omega t}\Phi_A)$ and equals $I$ at $t=0$. The
remaining identities substitute this into the definitions of $\Phi(T)$ and
$v$.
\end{proof}

\begin{algorithm}[H]
\caption{Fast HTF frequency sweep}\label{alg:sweep}
\begin{algorithmic}[1]
\Require Data, frequency grid $\{\omega_1,\dots,\omega_{N_\omega}\}$, output radius
$M_0$, tol $\eta$.
\Ensure $\Hcal(\omega_q)_{k\ell}$ for all $q$, $|k|,|\ell|\le M_0$.
\State \textbf{(Once)} Integrate $\dot\Phi_A=A(t)\Phi_A$, $\Phi_A(0)=I$, store
dense interpolant and $\Psi=\Phi_A(T,0)$.
\State \textbf{(Once)} Compute Floquet decomposition $\Psi=V\,\Lambda\,V^{-1}$
($\Lambda=\diag\mu_j$) for stable BVP solves.
\State \textbf{(Once)} For each $\ell$, precompute $g_\ell(s)=\Phi_A(s,0)^{-1}B(s)e^{i\ell\Omega s}e_c$
on the integration grid and its Fourier representation.
\For{$q=1,\dots,N_\omega$}
  \State Set $\Phi_T=e^{-i\omega_q T}\Psi$; solve $(I-\Phi_T)\xi_0=\Phi_T v(T)$
  via $V,\Lambda$ ($O(n^2)$ using $(I-e^{-i\omega_q T}\Lambda)^{-1}$).
  \State Get $v(T)=\widehat{g_\ell}(-\omega_q)$ from the precomputed Fourier data
  (interpolation / NUFFT).
  \State Form $\xi(t)=e^{-i\omega_q t}\Phi_A(t,0)(\xi_0+v(t))$ and
  $w(t)=C\xi+De^{i\ell\Omega t}e_c$; FFT to get $\Hcal(\omega_q)_{k\ell}$.
\EndFor
\State \Return all blocks.
\end{algorithmic}
\end{algorithm}

\begin{theorem}[Correctness of the sweep]\label{thm:sweep}
Algorithm~\ref{alg:sweep} returns, for each $\omega_q$ with
$e^{i\omega_q T}\notin\spec(\Psi)$, the same quantity as
Algorithm~\ref{alg:mono} at $\omega_q$, hence the error bound \eqref{eq:monoerr}
holds for every grid point with $\delta(\omega_q)^{-1}$ replaced by
$\min_j |1-e^{-i\omega_q T}\mu_j|^{-1}$.
\end{theorem}

\begin{proof}
By Proposition~\ref{prop:shift}, the only $\omega$-dependent objects in
Lemma~\ref{lem:monoform} are (a) the scalar phase $e^{-i\omega t}$,
(b) $\Phi_T=e^{-i\omega T}\Psi$, and (c) $v(T)=\widehat{g_\ell}(-\omega)$. Steps
5--7 reproduce \eqref{eq:xiform} verbatim with these substituted, so the
analytic output equals that of Algorithm~\ref{alg:mono}. The eigendecomposition
$\Psi=V\Lambda V^{-1}$ gives
$(I-\Phi_T)^{-1}=V(I-e^{-i\omega_q T}\Lambda)^{-1}V^{-1}$, a backward-stable
$O(n^2)$ solve once $V,\Lambda$ are computed; its conditioning is
$\min_j|1-e^{-i\omega_q T}\mu_j|^{-1}$ times the eigenvector conditioning
$\norm V\,\norm{V^{-1}}$. Substituting into Theorem~\ref{thm:mono} gives the
claim.
\end{proof}

\section{Complexity and comparison with the LTI case}\label{sec:complexity}

Let $n$ be the state dimension, $m,p$ input/output dimensions, $K$ the harmonic
bandwidth of the data, $M_0$ the requested output-harmonic radius,
$N_\omega$ the number of frequencies, and $N_F=O(M_0+K)$ the FFT length.

\begin{theorem}[Complexity]\label{thm:complexity}
\leavevmode
\begin{enumerate}
\item \textbf{Algorithm~\ref{alg:trunc} (truncation), per frequency.} With
window $M=M_0+O(\tfrac1\alpha\log\tfrac1\varepsilon)$ from \eqref{eq:Mchoice}: a
dense factorization costs $O((Mn)^3)$. For band-limited data ($\hat A_k=0$,
$|k|>K$) the matrix $W$ is banded with bandwidth $O(Kn)$, so the banded
$LU$ costs $O\!\big(M\,(Kn)^2\big)$ and the solve $O\!\big(M (Kn) (mM_0)\big)$.
Total over the grid: $O\!\big(N_\omega\,M\,(Kn)^2\big)$.
\item \textbf{Algorithm~\ref{alg:mono} (monodromy), per frequency.} One
$n\times n$ ODE solve with $N_s$ steps: $O(N_s n^3)$ (matrix products /
exponentials); $m$ quadratures $O(N_s n^2 m)$; $m$ BVP solves $O(n^3 m)$; FFTs
$O(p\,m\,N_F\log N_F)$. Dominant term $O(N_s n^3 + p m N_F\log N_F)$. Total over
the grid (no precomputation reuse): $O\!\big(N_\omega(N_s n^3+pmN_F\log N_F)\big)$.
\item \textbf{Algorithm~\ref{alg:sweep} (fast sweep).} \emph{One} ODE solve and
\emph{one} eigendecomposition of $\Psi$, $O(N_s n^3+n^3)$, amortized over the
grid; then per frequency only $O(n^2 m + pmN_F\log N_F)$. Total
\begin{equation}\label{eq:sweepcost}
O\!\big(\underbrace{N_s n^3}_{\text{once}}+N_\omega\,(n^2 m+p m N_F\log N_F)\big).
\end{equation}
\end{enumerate}
\end{theorem}

\begin{proof}
(1) The displayed matrix $W$ has size $(2M{+}1)n$; dense $LU$ is cubic, and
band-limited data make it banded with the stated bandwidth, for which banded
$LU$ and triangular solves have the quoted costs (standard banded linear
algebra). $M$ is set by Theorem~\ref{thm:trunc}. (2) Each line of
Algorithm~\ref{alg:mono} has the cost of the corresponding dense linear-algebra
/ FFT operation: the matrix ODE has $n^2$ entries and $N_s$ steps with $O(n^3)$
work per step (matrix multiply or Magnus exponential); the BVP solve is an
$n\times n$ system; FFTs are $O(N_F\log N_F)$ per output/input pair. (3) By
Proposition~\ref{prop:shift} the ODE and $g_\ell$ are computed once; the
per-frequency work is the $O(n^2)$ application of
$(I-e^{-i\omega_q T}\Lambda)^{-1}$ in the eigenbasis (diagonal solve plus two
$n\times n$ multiplies by $V,V^{-1}$), the $O(1)$-per-entry evaluation of
$v(T)=\widehat{g_\ell}(-\omega_q)$ (NUFFT/interpolation), and the FFTs. Summing
gives \eqref{eq:sweepcost}.
\end{proof}

\begin{corollary}[Parity with the LTI case]\label{cor:lti}
For an LTI system ($A,B,C,D$ constant, $K=0$) we have $\Phi_A(t,0)=e^{At}$,
$\Psi=e^{AT}$, the HTF is block-diagonal,
$\Hcal(\omega)_{kk}=C(i(\omega+k\Omega)I-A)^{-1}B+D$, and
Algorithm~\ref{alg:sweep} reduces to: one eigendecomposition of $A$ (once), then
per frequency a diagonal resolvent $O(n^2)$ evaluation. This is exactly the cost
of the classical pole–residue / state-space frequency-response evaluation for
LTI systems, $O(n^3+N_\omega n^2)$. Hence the LTP sweep \eqref{eq:sweepcost}
matches the LTI complexity up to the harmonic factor $N_F\log N_F$ (which is
$1$ when $M_0=K=0$), as required.
\end{corollary}

\begin{proof}
For $K=0$, the BVP \eqref{eq:bvp} has constant coefficients; its periodic
solution is $\xi=(i\omega I-A)^{-1}\hat b$ when $\ell=0$ and $\xi=0$ for
$\ell\ne0$ at the matching harmonic, recovering the diagonal LTI formula. The
$\omega$-update through $\Lambda=\spec(e^{AT})$ coincides with shifting the LTI
poles by $i\Omega\Z$, and the eigendecomposition is computed once. The stated
costs are then immediate from \eqref{eq:sweepcost}.
\end{proof}

\section{Summary of guarantees}

We have produced three algorithms and proved, for each, the four requested
properties.
\begin{itemize}
\item \textbf{Correctness / well-posedness.} The HTF \eqref{eq:HTFdef} is a
bounded operator under Assumption~\ref{ass:hyp} (Lemmas
\ref{lem:spectrum}--\ref{lem:bounded}) and equals the harmonic transform of a
finite periodic BVP (Proposition~\ref{prop:green}, Lemma~\ref{lem:monoform}),
which all three algorithms compute (Theorems
\ref{thm:trunc},\ref{thm:mono},\ref{thm:sweep}).
\item \textbf{Provable error bounds.} Truncation: exponential in the window
overhead, $C_1 e^{-\alpha(M-M_0)}$ (Theorem~\ref{thm:trunc}), with the rate
$\alpha$ tied to the generator gap $\gamma(\omega)$ via off-diagonal resolvent
decay (Lemma~\ref{lem:decay}, extended in Remark~\ref{rem:general}). Monodromy:
no harmonic-truncation error in $A$; total error linear in the ODE/quadrature
tolerance times $\delta(\omega)^{-1}$ plus controlled aliasing
(Theorem~\ref{thm:mono}).
\item \textbf{Avoiding the doubly-infinite operator.} Algorithms
\ref{alg:mono}--\ref{alg:sweep} never form $\Acal$ or its truncation; they
operate on $n\times n$ objects (Proposition~\ref{prop:green}).
\item \textbf{Structure exploitation.} Sparsity/band-limitedness gives banded
linear algebra in Algorithm~\ref{alg:trunc}; harmonic roll-off gives the
truncation rate and aliasing control; the monodromy/Floquet factorization
$\Psi=V\Lambda V^{-1}$ and the frequency-shift identity \eqref{eq:shift} give the
single-ODE fast sweep (Proposition~\ref{prop:shift},
Theorem~\ref{thm:complexity}).
\item \textbf{LTI-comparable efficiency.} The sweep costs
$O(N_s n^3 + N_\omega(n^2m + pmN_F\log N_F))$, reducing exactly to the classical
LTI frequency-response cost when $K=0$ (Corollary~\ref{cor:lti}).
\end{itemize}

\end{document}
